%%%%%%%%%%%%%%%%%%%%%%%%%%%%%%%%%%%
% Number Systems
%%%%%%%%%%%%%%%%%%%%%%%%%%%%%%%%%%%
\section{A Brief History of Number Systems}
%%%%%%%%%%%%%%%%%%%%%%%%%%%%%%%%%%%%%%%%%%%
\subsection{Early Development}
%%%%%%%%%%%%%%%%%%%%%%%%%%%%%%%%%%%%%%%%%%%
Early humans used tally marks to record quantities, such as notches on bones or stones, dating back to approximately 30,000 BCE \cite{ifrah2000universal}. More structured systems emerged with the Babylonians, who developed a base-60 system around 2000 BCE, and the Egyptians, who used hieroglyphic numerals around 3000 BCE \cite{ore1990number}. The base-12 system, or duodecimal system, emerged due to its practicality in trade, measurement, and division of goods. Its divisibility by 2, 3, 4, and 6 made it ideal for splitting quantities evenly \cite{wells2010number}. These systems laid the foundation for positional notation and arithmetic, with remnants of base-60 and base-12 still evident today in the division of time into 60 minutes and 60 seconds, the 12 months of a year, the 12 inches in a foot, and the grouping of items into dozens in commerce.

The Hindu-Arabic numeral system, developed between 500–1200 CE, introduced the concept of zero and a fully positional decimal system. This innovation transformed arithmetic and facilitated the development of more advanced mathematical concepts, forming the basis of the modern number system \cite{kline1990mathematics}.

The binary number system, foundational to digital computing, was first described by the Indian mathematician Pingala in the 2nd century BCE. Later, in the 17th century, Gottfried Wilhelm Leibniz formalized binary arithmetic, recognizing its potential for simplifying calculations \cite{kline1990mathematics}. This system’s alignment with the on-off states of electrical circuits made it the natural choice for modern computers.

%%%%%%%%%%%%%%%%%%%%%%%%%%%%%%%%%%%%%%%%%%%%
\subsection{Evolution of Mathematical Sets}
%%%%%%%%%%%%%%%%%%%%%%%%%%%%%%%%%%%%%%%%%%%%

The evolution of mathematical sets reflects humanity's growing understanding of quantity, ratios, and continuity.

\paragraph{Integers (\( \mathbb{Z} \)).}
The Babylonians and Egyptians used whole numbers for trade and astronomy. However, the concept of negative numbers was not widely accepted until Indian mathematician Brahmagupta introduced them in the 7th century CE to represent debts \cite{burton2010-number}. 

\paragraph{Rational Numbers (\( \mathbb{Q} \)).}
Rational numbers were extensively studied by the Greeks, particularly in their exploration of ratios in geometry. The discovery of irrational numbers, such as \( \sqrt{2} \), challenged their worldview, leading to the distinction between rational and irrational quantities \cite{maor2007}.

\paragraph{Real Numbers (\( \mathbb{R} \)).}
Real numbers evolved from the need to describe continuous quantities. Euclid's work on incommensurable lengths hinted at irrational numbers, but rigorous definitions came in the 19th century through the works of Dedekind (via Dedekind cuts) and Cantor (via set theory) \cite{rosen2018}.

\paragraph{Complex Numbers (\( \mathbb{C} \)).}
Complex numbers arose in the 16th century as mathematicians like Cardano and Bombelli sought solutions to polynomial equations, particularly cubic and quartic equations \cite{boyer1991history}. Initially dismissed as "imaginary," they were formalized in the 18th and 19th centuries by mathematicians such as Euler, Argand, and Gauss. Euler introduced the formula \( e^{i\theta} = \cos\theta + i\sin\theta \), Argand developed the geometric interpretation of complex numbers on the complex plane, and Gauss expanded their application to number theory and algebra \cite{kline1990mathematics, stewart2015-foundations}. Today, complex numbers are essential in fields like quantum mechanics, engineering, and signal processing \cite{maor2007}.


%%%%%%%%%%%%%%%%%%%%%%%%%%%%%%%%%%%%%%%%%%%%%%%%%%%%%%
\subsection{Pre-Computer Number Representations}
%%%%%%%%%%%%%%%%%%%%%%%%%%%%%%%%%%%%%%%%%%%%%%%%%%%%%%

The representation of numbers has evolved alongside advances in mathematics, science, and technology. While early systems like the Babylonian base-60 and Hindu-Arabic decimal notation provided the foundation for positional number systems, more specialized representations emerged to address the computational challenges of different eras. From manual calculations to the computer age, each system reflects the requirements and constraints of its time.

Before the advent of digital computers, numerical representations were primarily designed for manual or mechanical computation. These early methods focused on simplifying arithmetic while addressing the limitations of the tools available.

\paragraph{Decimal Positional Systems.}
The Hindu-Arabic numeral system, which became dominant during the medieval period, introduced positional notation and the concept of zero. This innovation enabled efficient manual computation and set the stage for future representations in mechanical and digital systems \cite{ifrah2000, burton2010-history}.

\paragraph{10's Complement for Subtraction.}
The concept of complements, such as 10's complement, predates the computer era and was used in mechanical calculators like the 19th-century Arithmometer. By representing negative numbers as the difference between a base (e.g., \( 10^n \)) and a positive value, subtraction could be performed as addition with a small adjustment. This method inspired later binary complement systems used in digital computers \cite{randell1982origins}.

\paragraph{Logarithmic Representations.}
In the early 17th century, John Napier introduced logarithms, enabling the multiplication of large numbers by adding their logarithms. Logarithmic tables and slide rules became essential tools for computation, representing a precursor to modern computational methods \cite{stewart2008-taming}.

%%%%%%%%%%%%%%%%%%%%%%%%%%%%%%%%%%%%%%%%%%%%
\subsection{Number Representations in Computing}
%%%%%%%%%%%%%%%%%%%%%%%%%%%%%%%%%%%%%%%%%%%%

The transition from mechanical to digital systems necessitated new numerical representations tailored to the constraints and capabilities of electronic computers.

\paragraph{Binary and Two's Complement.}
With the advent of binary systems, representing numbers using only two states (0 and 1) became standard in computing. Negative numbers posed a challenge, initially addressed with signed magnitude representation. However, this method proved inefficient for arithmetic operations, leading to the adoption of two's complement \cite{kline1990mathematics}. In two's complement, the negative of a number is represented as its binary complement plus one, simplifying subtraction and enabling efficient hardware implementation.

\paragraph{Excess-\( k \) Representations.}
First used in analog computing systems, excess-\( k \) representations became a cornerstone of digital computing with the IEEE 754 floating-point standard. By biasing the representation of numbers, excess-\( k \) allows for simplified comparisons and normalization, making it ideal for scientific computations \cite{goldberg1991}.

\paragraph{Fixed-Point Arithmetic.}
Fixed-point representation emerged as a practical solution for embedded systems and early computers with limited hardware capabilities. By dedicating a fixed number of bits to the fractional and integer parts of a number, fixed-point arithmetic provided a simpler alternative to floating-point for real-time applications \cite{smith1997science}.

\paragraph{Floating-Point Systems.}
The need for a dynamic range in numerical representation led to the development of floating-point systems. Introduced formally in the mid-20th century and standardized with IEEE 754 in 1985, floating-point representation encodes numbers as a sign, exponent, and mantissa. This system balances precision, range, and computational efficiency, making it indispensable for modern scientific computing \cite{goldberg1991}.

\paragraph{Binary Coded Decimal (BCD).}
BCD, introduced in the early computer era, represents decimal digits directly in binary form. This representation avoids rounding errors associated with binary floating-point arithmetic, making it ideal for financial and accounting systems \cite{randell1982origins}. Despite its storage inefficiency, BCD remains in use for applications requiring exact decimal precision.

%%%%%%%%%%%%%%%%%%%%%%%%%%%%%%%%%%%%%%
\subsection{Modern Trends and Challenges}
%%%%%%%%%%%%%%%%%%%%%%%%%%%%%%%%%%%%%%
Arbitrary precision and energy-efficient representations are current areas of continued research and development.

\begin{itemize}
    \item Arbitrary Precision Arithmetic: Modern software libraries support representations that exceed hardware constraints, enabling computations with arbitrarily high precision for cryptography and scientific research \cite{smith1997science}.
    \item Energy-Efficient Representations: In the era of low-power devices and AI accelerators, representations like fixed-point and logarithmic encoding are experiencing a resurgence due to their energy efficiency.
    \item Approximate Computing: To reduce power consumption and improve efficiency, approximate computing techniques trade precision for energy savings in tasks like image processing, machine learning, and data analytics. These methods employ representations that prioritize computational simplicity over exactness, reflecting the growing importance of resource-constrained systems \cite{han2013approximate}.
\end{itemize}